\documentclass[aps,prd,12pt,twocolumn,superscriptaddress,floatfix,longbibliography]{revtex4-2}

\usepackage{iftex}
\ifPDFTeX
  \usepackage{mathptmx}
\else
  \usepackage{fontspec}
  \setmainfont{TeX Gyre Termes}
  \setsansfont{TeX Gyre Heros}
  \setmonofont{Latin Modern Mono}
\fi
\usepackage{amsmath}
\usepackage{booktabs}
\usepackage{graphicx}
\usepackage{xurl}
\usepackage[a4paper, margin=0.6in]{geometry}

\setlength{\emergencystretch}{2em}
\hbadness=10000
\renewcommand{\footnoterule}{}
\graphicspath{{artifacts/figures/predecoder_v2/}{artifacts/figures/predecoder/}{figure/}}
\providecommand{\href}[2]{#2}

\newcommand{\predecoderfigure}[2]{%
\IfFileExists{artifacts/figures/predecoder_v2/#1_v2.pdf}{%
    \includegraphics[width=#2]{artifacts/figures/predecoder_v2/#1_v2.pdf}%
}{%
    \IfFileExists{artifacts/figures/predecoder_v2/#1_v2.png}{%
        \includegraphics[width=#2]{artifacts/figures/predecoder_v2/#1_v2.png}%
    }{%
        \IfFileExists{artifacts/figures/predecoder/#1.pdf}{%
            \includegraphics[width=#2]{artifacts/figures/predecoder/#1.pdf}%
        }{%
            \IfFileExists{artifacts/figures/predecoder/#1.png}{%
                \includegraphics[width=#2]{artifacts/figures/predecoder/#1.png}%
            }{%
                \fbox{\parbox{0.92\linewidth}{%
                \centering Figure file missing:\\
                \texttt{artifacts/figures/predecoder_v2/#1_v2.png}}}%
            }%
        }%
    }%
}%
}

\begin{document}

\title{Design and Evaluation of a Transition-Aware Neural Pre-Decoder for Surface-Code Quantum Error Correction}
\author{Jaemin Jeon}
\affiliation{Korea University}
\date{May 2026}

\begin{abstract}
PyMatching is a standard matching-based decoder for estimating logical
information from repeated surface-code syndrome measurements. This work is
motivated by an empirical observation: for some shots where PyMatching fails,
flipping one or two local detector events before matching can move the logical
decision back to the correct class. The same kind of edit can also damage an
already-correct shot, so the problem is separated into local edit generation
and edit adoption.

The proposed transition-aware pre-decoder builds small detector-edit candidates
from a 36-channel syndrome/noise volume and scores each candidate by the change
it induces in the downstream PyMatching logical class. Here, transition-aware
does not mean a separate temporal transition model; it means that candidate
selection uses the downstream decoder decision before and after a local edit.
The neural model does not output a correction chain or a final logical frame.
It only passes either an edited syndrome or the raw no-edit syndrome to
PyMatching, which remains the final decoder. Adoption thresholds are fixed
using train/validation families only.

On the held-out \texttt{stage\_c\_corr} evaluation, selected-mode pre-decoding
improves mean accuracy over raw PyMatching by \(+0.66\) pp for d3 and
\(+0.56\) pp for d5. The d5 result is a conservative selected-mode result:
only two of eight seeds adopt local edits. At d7, the selected-mode gain is only
\(+0.02\) pp. Oracle analysis indicates that recoverable local edits still
exist at d7. The current model often generates useful candidates, but the
selector does not choose them reliably under held-out noise-family shift.
\end{abstract}

\maketitle

\section{Introduction}

Fault-tolerant quantum computation requires a decoder that can infer the
logical effect of physical errors from noisy stabilizer measurements. In a
surface code, local stabilizers are repeatedly measured, and changes in the
syndrome define detector events in a space-time decoding volume
~\cite{surfacecode}. These detector events do not directly reveal the physical
error path; rather, they constrain the set of error chains that are compatible
with the observed syndrome.

Matching-based decoders, including PyMatching, are standard tools for this
inference problem~\cite{pymatching}. The narrower question is whether a neural
module can make a small local edit to the detector-event input before
matching, while leaving the final logical-frame decision to PyMatching.

This viewpoint follows neural pre-decoding, where a front-end module modifies
or reduces a syndrome before a downstream global decoder is applied. NVIDIA's
AI-based pre-decoder work demonstrates this idea at a much larger surface-code
scale~\cite{nvidiaising}. Here, the focus is on a smaller setting where local
edit generation and edit adoption can be analyzed separately: detector-edit
candidates are generated, scored, adopted or rejected by a safety policy, and
evaluated through the downstream PyMatching outcome.

The implementation combines a 3D residual syndrome trunk, local motif candidate
generation, and a patch-head selector. The selector target is defined by
downstream PyMatching benefit or harm signals. The analysis also separates raw,
candidate, selected, and oracle branches so that gains and failure modes can be
attributed more precisely. This separation is important for d7, where useful
candidates exist but the selector does not choose them reliably.

\section{Background and Related Work}

\subsection{Surface Codes and Matching Decoders}

A surface code repeatedly measures stabilizers on a two-dimensional qubit
lattice. A detector event is produced when the measurement outcome associated
with the same stabilizer changes across time; detector definitions also include
events closed by boundaries and final measurements. A decoder uses these
events, together with the code boundaries, to infer a likely error chain and
the logical frame induced by that chain.

PyMatching formulates this problem as a minimum-weight matching problem on a
sparse matching graph~\cite{pymatching}. Here, PyMatching plays three roles. It
is the raw no-edit baseline, the evaluation criterion used to assign
benefit/harm labels to candidate edits, and the downstream decoder that
evaluates the edited syndrome selected by the neural pre-decoder. Consequently,
reported performance refers to the hybrid pipeline composed of a neural front
end and a matching decoder.
In this work, \texttt{logical\_class4} denotes the four-class logical
decision label used for evaluation. The text also refers to this quantity as
the logical frame or logical decision when discussing the PyMatching output.

\subsection{Neural Decoding and Pre-Decoding}

Neural approaches to surface-code decoding can be divided into two broad
families. The first family attempts to predict the logical class or correction
field directly from syndrome data. AlphaQubit and AlphaQubit 2 are
representative large-scale examples of this replacement-style neural decoding
direction~\cite{alphaqubit,alphaqubit2}. Convolutional neural decoders have
also been studied as lattice-aware direct decoders for surface
codes~\cite{convnnsurface}. The second family keeps a global decoder in the
loop and uses a neural model as a pre-decoder that modifies the syndrome or
reduces the residual decoding problem. NVIDIA Ising-Decoding is the closest
structural reference for this second direction~\cite{nvidiaising}.

The proposed approach belongs to the second family. A direct neural classifier can be
tested at small distances, but it is difficult to make it a stable reference as
distance and noise-family shift increase. A pre-decoder preserves the
stability of PyMatching while allowing targeted edits on local syndrome
patterns. Here, the neural model is restricted to local candidate selection.

Table~\ref{tab:related-work} summarizes the position of the proposed approach relative to
related studies. The table is not a head-to-head performance comparison on the
same dataset. It clarifies which decision is learned and which decoder produces
the final logical decision. The numerical claims below are limited to the
Stage A/B/C noise families and the held-out \texttt{stage\_c\_corr} evaluation.

\begin{table*}[!tbp]
\caption{Position of this work relative to related decoding studies.}
\label{tab:related-work}
\centering
\scriptsize
\begin{tabular*}{\textwidth}{@{\extracolsep{\fill}}llll}
\toprule
work & learned decision & final decoder & relation \\
\midrule
AlphaQubit family~\cite{alphaqubit,alphaqubit2} & logical/correction decision & neural decoder & direct final-decision model \\
CNN surface decoder~\cite{convnnsurface} & local-pattern correction & CNN decoder & direct decoder baseline \\
NVIDIA Ising-Decoding~\cite{nvidiaising} & pre-decoding decision & residual decoder & pre-decoder before residual decoder \\
Near-term neural decoder~\cite{neartermnn} & experimental syndrome mapping & neural decoder & experimental syndrome-to-logical mapping \\
Google Willow~\cite{willowthreshold} & hardware-scale QEC evidence & hardware decoder stack & hardware-scale motivation only \\
This work & local detector edit & PyMatching & local-edit selection before PyMatching \\
\bottomrule
\end{tabular*}
\end{table*}

\section{Dataset and Noise Setting}

The experiments use rotated surface-code memory circuits generated with Stim
and custom staged noise families~\cite{stim}. The noise families are built in
three stages. \texttt{stage\_a\_si1000} is the baseline SI1000-style physical
noise setting, following the superconducting-inspired setting used in
surface-code resource estimates~\cite{honeycombmemory}. \texttt{stage\_b\_local} adds local heterogeneity, and
\texttt{stage\_c\_corr} adds a correlated stray surrogate. The quantitative
claims in this work are based on held-out \texttt{stage\_c\_corr}
evaluation.

Each shot is represented as a 36-channel syndrome/noise volume containing
detector events, a valid-detector mask, geometry metadata, noise-family
indicators, distance/round statistics, event fractions, and physical-noise
summary channels. Channel semantics are fixed across code distances, while the
spatial and temporal volume sizes grow with distance. The input shapes for d3,
d5, and d7 are \([36,4,4,4]\), \([36,6,6,6]\), and \([36,8,8,8]\), respectively.

Each original manifest contains \(2048\) shots per family. Training local-edit
targets are generated from fixed \(1024\)-shot subsets of the train/validation
families. The \(1024\) held-out shots per seed in \texttt{stage\_c\_corr} are
used only for final evaluation and post hoc oracle analysis. Each train family
uses a fixed validation split of \(154\) shots, and the remaining samples are
used to train the trunk and selector.

\begin{table*}[!tbp]
\caption{Dataset configuration by code distance. Each original dataset contains
\(2048\) shots per family. Training local-edit targets are generated from fixed
\(1024\)-shot subsets of the train/validation families, while the held-out
\(1024\) shots are used only for final evaluation and oracle analysis.}
\label{tab:dataset-config}
\centering
\scriptsize
\begin{tabular*}{\textwidth}{@{\extracolsep{\fill}}lccccc}
\toprule
distance & rounds & input shape & train/validation families & held-out family & reported seeds \\
\midrule
d3 & 3 & \([36,4,4,4]\) & A/B, \(1024\) shots each & C, \(1024\) shots & \(0..7\) \\
d5 & 5 & \([36,6,6,6]\) & A/B, \(1024\) shots each & C, \(1024\) shots & \(0..7\) \\
d7 & 7 & \([36,8,8,8]\) & A/B, \(1024\) shots each & C, \(1024\) shots & \(0..57\) \\
\bottomrule
\end{tabular*}
\end{table*}

\begin{table*}[!tbp]
\caption{Noise families used in the experiments.}
\label{tab:noise-family}
\centering
\footnotesize
\begin{tabular*}{\textwidth}{@{\extracolsep{\fill}}llll}
\toprule
family & noise structure & role & reported use \\
\midrule
\texttt{stage\_a\_si1000} & SI1000-style baseline physical noise & baseline family & train/validation \\
\texttt{stage\_b\_local} & Stage A plus local heterogeneity & structured-noise family & train/validation \\
\texttt{stage\_c\_corr} & Stage B plus correlated stray surrogate & primary test family & held-out evaluation \\
\bottomrule
\end{tabular*}
\end{table*}

\section{Proposed Method}

The proposed method is a local-edit neural pre-decoder placed before
PyMatching. The full inference path is shown in
Figure~\ref{fig:pipeline}. The neural model generates small local edits from a
detector-event volume and evaluates whether each candidate is beneficial or
harmful for downstream PyMatching. Only candidates that pass the safety policy
are applied to the syndrome; otherwise, the raw no-edit syndrome is passed to
PyMatching.
The selector uses the transition between the raw PyMatching logical class and
the logical class after a local edit as part of both its training signal and
candidate features.

\begin{figure*}[!tbp]
\centering
\predecoderfigure{fig1_predecoder_pipeline}{\textwidth}
\caption{Pipeline of the proposed neural pre-decoder and PyMatching. The neural
module only decides whether to apply a local syndrome edit; PyMatching remains
the final logical-frame decoder.}
\label{fig:pipeline}
\end{figure*}

\subsection{SyndromeEditPreDecoder}

\texttt{SyndromeEditPreDecoder} consists of a 3D convolution stem followed by
three residual 3D convolution blocks. The final d3/d5 setting uses hidden
width \(24\), dense hidden dimension \(64\), dropout \(0.1\), and \(118834\)
parameters. The trunk produces detector-level \texttt{edit\_logits},
shot-level \texttt{needs\_edit\_logits}, and pooled shot features. The pooled
features and local candidate features are then passed to the patch-head
selector.

\subsection{Local Motif Candidates}

The model restricts syndrome edits to small local changes. The final setting
uses \texttt{local\_motif\_selector}, which builds a restricted set of local
motif candidates derived from train/validation hard-shot edit targets. The
identity candidate, corresponding to no edit, is always included, so raw
PyMatching remains the internal fallback path of the selected mode.
During target construction, local search is applied to raw PyMatching failures
in the train/validation families with temporal radius \(1\), spatial radius
\(1\), edit weight at most \(2\), and at most \(20\) variants per incorrect
shot. A separate held-out search on \texttt{stage\_c\_corr} found that local
edits could recover 65 of 73 raw d3 failures, 92 of 114 raw d5 failures, and
114 of 130 raw d7 failures at the oracle-analysis level. These held-out search
results are not used for training, motif selection, or threshold selection;
they are used only as post hoc evidence that recoverable local edits exist.

\subsection{Benefit/Harm Selector}

\texttt{CandidateEditSelector} assigns a scalar score to each local edit
candidate. A candidate is labeled as useful when applying the edit improves the
final \texttt{logical\_class4} decision after PyMatching relative to raw
no-edit PyMatching. The selector input is the concatenation of a pooled shot feature
from the 3D trunk and candidate-level features. The candidate features include
edit weight, summaries of detector edit probabilities, motif frequency,
coordinate summaries, nearby event/evidence values, and a small
\(3\times3\times3\) local patch embedding. The selector is trained with Adam
using learning rate \(5\times10^{-4}\) for at most \(6\) epochs. Its loss
combines group cross-entropy for choosing the best candidate with a pairwise
benefit/harm ranking loss using margin \(0.5\). Harmful candidates are weighted
more strongly than missed beneficial candidates in the selector loss. This
makes the selector a decoder-aware ranking module. Figure~\ref{fig:architecture}
summarizes the model structure.

\begin{figure*}[!tbp]
\centering
\predecoderfigure{fig2_model_architecture}{\textwidth}
\caption{Architecture of the neural pre-decoder. The
\texttt{SyndromeEditPreDecoder} builds a shared 3D feature volume and emits
edit, needs-edit, and pooled-shot heads. The patch-head
\texttt{CandidateEditSelector} then scores candidate features before the
selected edit or fallback path is passed to PyMatching.}
\label{fig:architecture}
\end{figure*}

\subsection{Selected-Mode Safety}

The selected mode applies a candidate only after it passes the adoption rule.
Adoption thresholds are fixed using train/validation families only; held-out
\texttt{stage\_c\_corr} results are not used to choose thresholds or seed-level
adoption decisions. If validation gain over raw no-edit is insufficient, or if
harm and selector-margin conditions fail, the selected mode falls back to raw
no-edit. This safety rule reflects the fact that local edits can help some
seeds while harming others.

\section{Experimental Setup}

Four branches are compared. Raw PyMatching decodes the original syndrome.
Selected predecoder is the final proposed output, which applies either the
selected local edit or raw no-edit according to the safety policy. Candidate
branch applies the learned candidate without fallback. Target local-edit oracle
is the best achievable local edit inside the fixed candidate space.

The primary metric is \texttt{logical\_class4} accuracy on held-out
\texttt{stage\_c\_corr}. Candidate and selected branches are reported
separately because they answer different questions: the candidate branch
measures the potential and risk of the learned candidate space, whereas the
selected branch measures the behavior of the full safety-gated system.

\section{Results}

\subsection{Main Accuracy}

Table~\ref{tab:main-results} and Figure~\ref{fig:main-accuracy} show the main
results for d3, d5, and d7. At d3 and d5, the selected predecoder improves mean
accuracy over raw PyMatching. At d7, the selected-mode gain is
\(+0.02\) pp, well below what would count as robust learned recovery.

\begin{table*}[!tbp]
\caption{Main accuracy on held-out \texttt{stage\_c\_corr}.}
\label{tab:main-results}
\centering
\scriptsize
\begin{tabular*}{\textwidth}{@{\extracolsep{\fill}}lcccccc}
\toprule
distance & seeds & raw & selected & candidate & oracle & gain \\
\midrule
d3 & \(0..7\) & 92.87\% & 93.53\% & 93.53\% & 99.22\% & \(+0.66\) pp \\
d5 & \(0..7\) & 88.87\% & 89.43\% & 89.18\% & 97.85\% & \(+0.56\) pp \\
d7 & \(0..57\) & 87.30\% & 87.32\% & 87.15\% & 98.44\% & \(+0.02\) pp \\
\bottomrule
\end{tabular*}
\end{table*}

\begin{figure*}[!tbp]
\centering
\predecoderfigure{fig3_main_accuracy_comparison}{\textwidth}
\caption{Held-out accuracy by code distance. The selected predecoder improves
over raw PyMatching at d3 and d5, while d7 remains close to raw even though
oracle analysis finds recoverable local edits.}
\label{fig:main-accuracy}
\end{figure*}

\subsection{D3/D5 Robustness}

Here, seed classes denote positive/neutral/negative selected-delta seed counts.
For d3, all seeds \(0..7\) have positive selected deltas even though the
shot-level outcomes include both improved and harmed shots (\(205/151\)); the
net seed-level effect is positive throughout. For d5, only two seeds adopt local
selector edits, while the other six seeds fall back to raw no-edit. The d5 mean
gain is therefore a selected-mode result from the adopted seeds, not uniformly
active learned editing across all seeds.

\begin{table}[!tbp]
\caption{D3/D5 seed \(0..7\) robustness summary. Seed classes are
positive/neutral/negative seed counts.}
\label{tab:d3d5-robustness}
\centering
\footnotesize
\begin{tabular*}{\columnwidth}{@{\extracolsep{\fill}}lccc}
\toprule
distance & seed classes & gain & improved/harmed \\
\midrule
d3 & \(8/0/0\) & \(+0.66\) pp & \(205/151\) \\
d5 & \(2/6/0\) & \(+0.56\) pp & \(56/10\) \\
\bottomrule
\end{tabular*}
\end{table}

The seed-level bootstrap 95\% confidence intervals are approximately
\([+0.45,+0.85]\) pp for d3 and \([0.00,+1.37]\) pp for d5. In the exact sign
test, the one-sided p-value is \(0.0039\) for d3 and \(0.25\) for d5. Thus, the
d3 result is consistently positive across the checked seeds. In contrast, the
d5 result is mainly a safety-gated selected-mode gain.

\begin{table}[!tbp]
\caption{D3/D5 paired seed-level statistics.}
\label{tab:paired-stats}
\centering
\footnotesize
\begin{tabular*}{\columnwidth}{@{\extracolsep{\fill}}lccc}
\toprule
distance & seed classes & sign p & sign-flip p \\
\midrule
d3 & \(8/0/0\) & 0.0039 & 0.0039 \\
d5 & \(2/6/0\) & 0.2500 & 0.2500 \\
\bottomrule
\end{tabular*}
\end{table}

\begin{table*}[!tbp]
\caption{D5 seed-level selected-mode fallback behavior.}
\label{tab:d5-fallback}
\centering
\scriptsize
\begin{tabular*}{\textwidth}{@{\extracolsep{\fill}}clrrl}
\toprule
seed & selected mode & selected delta & candidate delta & interpretation \\
\midrule
0 & raw no-edit & \(+0.00\) pp & \(+0.00\) pp & fallback \\
1 & raw no-edit & \(+0.00\) pp & \(+0.00\) pp & fallback \\
2 & local selector & \(+2.15\) pp & \(+2.15\) pp & useful edit \\
3 & local selector & \(+2.34\) pp & \(+2.34\) pp & useful edit \\
4 & raw no-edit & \(+0.00\) pp & \(-0.78\) pp & harmful blocked \\
5 & raw no-edit & \(+0.00\) pp & \(+0.00\) pp & fallback \\
6 & raw no-edit & \(+0.00\) pp & \(-1.17\) pp & harmful blocked \\
7 & raw no-edit & \(+0.00\) pp & \(+0.00\) pp & fallback \\
\bottomrule
\end{tabular*}
\end{table*}

Table~\ref{tab:d5-fallback} shows that the d5 mean gain comes from seeds 2 and
3. Seeds 4 and 6 have harmful candidate branches, but selected mode blocks
them by falling back to raw no-edit.

\subsection{Candidate Branch and Selected Branch}

The candidate branch reveals the potential of the candidate space, but applying
it without fallback also admits harmful edits. At d7, the candidate branch has
\(161/251\) improved/harmed shots, meaning that harmful effects dominate. The
selected mode avoids large degradation by choosing no-edit for most seeds, but
it also misses most recoverable edits found by the oracle analysis. The results
below therefore keep the candidate, selected, and oracle branches separate.
At d3, candidate and selected accuracy are both 93.53\%, so fallback does not
change the result. In contrast, selected accuracy exceeds candidate accuracy by
\(+0.25\) pp at d5 and \(+0.17\) pp at d7, quantifying the effect of the safety
policy in blocking harmful candidates.

\newpage

\section{D7 Limitation Analysis}

The d7 selected-mode delta is only \(+0.02\) pp. However, the target
local-edit oracle accuracy is \(98.44\%\), and all 58 d7 seeds have
positive oracle recoverability. Thus the d7 result is not simply explained by
the absence of useful local edits in the candidate space.

The more direct issue is that the selector does not reliably distinguish
beneficial candidates from harmful candidates on held-out noise. The d7
candidate outcomes split into \(6\) positive, \(35\) neutral, and \(17\)
harmful seeds, with mean actual candidate delta about \(-0.15\) pp.
Figure~\ref{fig:d7-oracle} summarizes the oracle recoverability and false-positive
structure.

\begin{figure*}[!tbp]
\centering
\predecoderfigure{fig4_d7_oracle_gap_false_positive}{\textwidth}
\caption{D7 oracle recoverability and false-positive candidate problem. Useful edits remain
inside the candidate space, but the learned selector does not reliably choose
validation-positive branches on held-out data.}
\label{fig:d7-oracle}
\end{figure*}

The validation-to-held-out mismatch shows the same issue at the seed level. Among 22
validation-positive candidate seeds, 13 are harmful on held-out evaluation and
only 5 remain positive. Thus validation evidence alone is insufficient for safe
candidate adoption. These 13 validation-positive but held-out harmful seeds are
referred to as harmful false positives; the 4 held-out neutral seeds are not
included in this harmful ratio. The harmful held-out outcome rate among
validation-positive seeds is \(13/22=59.09\%\).

\begin{table}[!tbp]
\caption{Mismatch between d7 validation class and held-out outcome.}
\label{tab:d7-mismatch}
\centering
\footnotesize
\begin{tabular*}{\columnwidth}{@{\extracolsep{\fill}}lccc}
\toprule
validation class & harmful & neutral & positive \\
\midrule
neutral & 4 & 31 & 1 \\
positive & 13 & 4 & 5 \\
\bottomrule
\end{tabular*}
\end{table}

\begin{figure*}[!tbp]
\centering
\predecoderfigure{fig5_d7_validation_heldout_scatter}{\textwidth}
\caption{Seed-level scatter of d7 validation delta versus held-out candidate
delta. Validation-positive evidence does not reliably transfer to held-out
correlated noise.}
\label{fig:d7-scatter}
\end{figure*}

There are 17 harmful candidate seeds, and the selected-mode guard blocks all
\(17/17\) of them by falling back to raw no-edit. The safety policy leaves d7
recovery unresolved, while preventing the harmful candidate branch from being
emitted as the final output.

Additional d7 threshold sweeps did not resolve this within the current model
family. A scalar adoption-threshold grid checked \(183040\)
policies, but none satisfied the predefined preserve/recover/block conditions.
The exact sentinel seed sets are recorded in the reproducibility package.
Cross-family hard positive-vs-negative objectives and candidate-compatibility
pairwise top-k objectives also failed to preserve true positives while blocking
false positives. The next d7 improvement requires a redesigned selector
ranking/generalization objective rather than simple threshold tuning.

\begin{figure*}[!tbp]
\centering
\predecoderfigure{fig6_oracle_recovery_distribution}{\textwidth}
\caption{Seed-level oracle recovery. D3 and d5 convert part of the recoverable
local-edit space into selected-mode gains, while d7 selected mode uses almost
none of the recoverable candidate space.}
\label{fig:oracle-recovery}
\end{figure*}

\newpage

\section{Discussion}

The d3/d5 results depend on the adoption rule as much as on edit generation.
At d5, the candidate branch contains harmful seeds, but the final selected mode
applies edits only to seeds 2 and 3, which satisfy the validation-side adoption
condition. The result supports a limited claim: within this candidate set, the
selector can identify cases where a small edit helps PyMatching.

The same procedure gives almost no gain at d7. Oracle analysis shows that
recoverable local edits still exist, but many validation-positive candidates
become harmful on the held-out family; the selector does not generalize well
enough to choose useful edits reliably across noise families.

The quantitative claims are limited to the Stage A/B/C noise families used in
this study. This range was chosen to separate candidate generation from
candidate adoption; hardware-scale decoder performance is outside the scope of
these experiments.

\section{Conclusion}

The proposed pre-decoder selects small detector edits before PyMatching is run.
On held-out \texttt{stage\_c\_corr}, selected-mode
pre-decoding improves mean accuracy over raw PyMatching by \(+0.66\) pp at d3
and \(+0.56\) pp at d5. As the seed-level analysis shows, the d5 gain remains a
selected-mode improvement from two adopted seeds.

At d7, the selected-mode gain is only \(+0.02\) pp. Recoverable edits exist in
the candidate space, but validation-positive candidates often become harmful
on held-out noise. The next improvement is to focus on selector targets and
validation-to-held-out generalization rather than simply increasing the
candidate set.

\section*{Reproducibility Note}

The result tables are based on archived pre-decoder summaries in the
repository: final result tables, paired statistics, oracle-recovery
distribution, d7 harmful-edit taxonomy, and the reproducibility package. The
final consistency check reports \(37/37\) checks and \(0\) failures.
Hyperparameter-sensitivity notes are kept as a separate supplementary memo.

\newpage

\bibliographystyle{unsrtnat}
\bibliography{main}

\end{document}
